\HLASourcePage{424}{423}
\thispagestyle{hla-body}
\begin{hlafixedpage}
  \HLAFixedBlock{90.000}{68.648}{514.000}{\rmfamily\raggedright}{13.98}{15.66}{%
    \textbf{Annex E}%
  }
  \HLAFixedBlock{90.000}{97.140}{514.000}{\rmfamily\raggedright}{12.00}{13.44}{%
    (normative) %
  }
  \HLAFixedBlock{90.000}{128.648}{514.000}{\rmfamily\raggedright}{13.98}{15.66}{%
    \textbf{Rationale}%
  }
  \HLAFixedBlock{90.000}{169.140}{514.000}{\rmfamily\raggedright}{12.00}{13.44}{%
    \textbf{E.1 Overview}%
  }
  \HLAFixedBlock{90.000}{195.786}{514.000}{\rmfamily\raggedright}{9.96}{11.16}{%
    This annex contains additional information that might be useful to explain various aspects of HLA and is not%
  }
  \HLAFixedBlock{90.000}{207.785}{514.000}{\rmfamily\raggedright}{9.96}{11.16}{%
    a complete rationale of the HLA. Nothing in this clause is required in an HLA implementation.%
  }
  \HLAFixedBlock{90.000}{231.785}{514.000}{\rmfamily\raggedright}{9.96}{11.16}{%
    This annex is organized to match the organization of the main body of this specification, and the subclause%
  }
  \HLAFixedBlock{90.000}{243.785}{514.000}{\rmfamily\raggedright}{9.96}{11.16}{%
    numbers in this annex correspond to the subclause numbers of the main body. Note, however, that not all%
  }
  \HLAFixedBlock{90.000}{255.785}{514.000}{\rmfamily\raggedright}{9.96}{11.16}{%
    subclauses in the main body have corresponding subclauses in this annex. In other words, text is present in%
  }
  \HLAFixedBlock{90.000}{267.785}{514.000}{\rmfamily\raggedright}{9.96}{11.16}{%
    this annex only if additional comments were deemed helpful to understanding the corresponding subclauses%
  }
  \HLAFixedBlock{90.000}{279.784}{514.000}{\rmfamily\raggedright}{9.96}{11.16}{%
    in the main body.%
  }
  \HLAFixedBlock{90.000}{309.383}{514.000}{\rmfamily\raggedright}{10.98}{12.30}{%
    \textbf{E.1.4 Background}%
  }
  \HLAFixedBlock{90.000}{334.806}{514.000}{\rmfamily\raggedright}{9.96}{11.16}{%
    Consideration was given to providing several iterating functions; however, there were found to be serious%
  }
  \HLAFixedBlock{90.000}{346.805}{514.000}{\rmfamily\raggedright}{9.96}{11.16}{%
    problems documenting and describing the iterating-based semantics of these services. Additionally, there%
  }
  \HLAFixedBlock{90.000}{358.805}{514.000}{\rmfamily\raggedright}{9.96}{11.16}{%
    was a problem with every redaction produced. The candidate definitions of how exceptions were handled%
  }
  \HLAFixedBlock{90.000}{370.805}{514.000}{\rmfamily\raggedright}{9.96}{11.16}{%
    and recovered from during the “middle” of an iteration sequence never satisfied a large enough set of users%
  }
  \HLAFixedBlock{90.000}{382.805}{514.000}{\rmfamily\raggedright}{9.96}{11.16}{%
    to be acceptable for retention in the specification. Each of the users wanted the behavior to meet their%
  }
  \HLAFixedBlock{90.000}{394.805}{514.000}{\rmfamily\raggedright}{9.96}{11.16}{%
    specific application needs, which in most cases were inconsistent and impossible to meet with a common%
  }
  \HLAFixedBlock{90.000}{406.805}{514.000}{\rmfamily\raggedright}{9.96}{11.16}{%
    mechanism. As a result, it seemed to make the most sense not to have any iterating services and to let users%
  }
  \HLAFixedBlock{90.000}{418.804}{514.000}{\rmfamily\raggedright}{9.96}{11.16}{%
    produce their own aggregate functionality with the exception response semantics they needed.%
  }
  \HLAFixedBlock{90.000}{454.140}{514.000}{\rmfamily\raggedright}{12.00}{13.44}{%
    \textbf{E.4 Federation management}%
  }
  \HLAFixedBlock{90.000}{486.383}{514.000}{\rmfamily\raggedright}{10.98}{12.30}{%
    \textbf{E.4.6 }%
  }
  \HLAFixedBlock{118.680}{486.383}{485.320}{\rmfamily\raggedright}{10.98}{12.30}{%
    \textbf{\emph{Destroy Federation Execution}}%
  }
  \HLAFixedBlock{274.379}{486.383}{329.621}{\rmfamily\raggedright}{10.98}{12.30}{%
    \textbf{ service}%
  }
  \HLAFixedBlock{90.000}{511.806}{514.000}{\rmfamily\raggedright}{9.96}{11.16}{%
    Since there can be no joined federates for the destroy service invocation to succeed, there might be a%
  }
  \HLAFixedBlock{90.000}{523.805}{514.000}{\rmfamily\raggedright}{9.96}{11.16}{%
    problem if the federation execution has “zombie” or otherwise uncooperative joined federates. To aid%
  }
  \HLAFixedBlock{90.000}{535.805}{514.000}{\rmfamily\raggedright}{9.96}{11.16}{%
    matters in this situation, a joined federate can be resigned by another joined federate via a MOM interaction.%
  }
  \HLAFixedBlock{90.000}{565.403}{514.000}{\rmfamily\raggedright}{10.98}{12.30}{%
    \textbf{E.4.31 }%
  }
  \HLAFixedBlock{124.800}{565.403}{479.200}{\rmfamily\raggedright}{10.98}{12.30}{%
    \textbf{\emph{Query Federation Restore Status}}%
  }
  \HLAFixedBlock{296.341}{565.403}{307.659}{\rmfamily\raggedright}{10.98}{12.30}{%
    \textbf{ service}%
  }
  \HLAFixedBlock{90.000}{590.826}{514.000}{\rmfamily\raggedright}{9.96}{11.16}{%
    One way to employ federation restore status is for a federation manager federate to not invoke the %
  }
  \HLAFixedBlock{491.820}{590.826}{112.180}{\rmfamily\raggedright}{9.96}{11.16}{%
    \emph{Restore}%
  }
  \HLAFixedBlock{90.000}{602.825}{514.000}{\rmfamily\raggedright}{9.96}{11.16}{%
    \emph{Complete}%
  }
  \HLAFixedBlock{128.280}{602.825}{475.720}{\rmfamily\raggedright}{9.96}{11.16}{%
     service until all the other joined federates have successfully completed their restore. In this way, a%
  }
  \HLAFixedBlock{90.001}{614.825}{513.999}{\rmfamily\raggedright}{9.96}{11.16}{%
    federation manager federate can “throttle” the restore process as it monitors the restore status.%
  }
\end{hlafixedpage}
\clearpage
